% Partie : méthodes
% Chapitre : démonstrations
% Section : commencer une démonstration

\documentclass[../../../include/open-logic-section]{subfiles}

\begin{document}

\olfileid[fr]{mth}{prf}{str}

\olsection{Commencer une démonstration}

Mais par où commencer ?

On vous a donné un énoncé à démontrer :
il doit donc être la dernière chose
établie dans votre démonstration.
Vous pouvez bien sûr \emph{annoncer}
au début que vous allez le démontrer,
mais vous ne devez pas l'utiliser
comme hypothèse. Écrivez ce que
vous cherchez à démontrer au bas
d'une feuille vierge ; vous garderez
ainsi votre objectif en vue.

Vous disposez peut-être ensuite
de certaines hypothèses. Leur rôle
deviendra plus clair dans la section
suivante, lorsque nous parlerons du
\emph{type} de démonstration à construire.
Écrivez-les en haut de la page et
signalez qu'il s'agit d'hypothèses :
si vous supposez $p$, écrivez par
exemple «~Supposons $p$~» ou
«~Admettons $p$~».
Enfin, l'énoncé peut faire intervenir
des définitions que vous devez
connaître. On peut vous demander
d'employer une définition précise,
ou plusieurs définitions peuvent
figurer dans les hypothèses ou
dans la conclusion visée.
\emph{Notez-les et assurez-vous
de comprendre ce qu'elles signifient.}

La forme logique de l'énoncé guide l'organisation de votre
démonstration. La section suivante montre comment commencer
le raisonnement dans chaque cas.

\end{document}
